\renewcommand{\abstractname}{Resum}

\begin{abstract}
Aquest Treball de Fi de Grau aborda les ineficiències derivades del procés manual de gestió de despeses dels empleats d'una empresa de consultoria informàtica. Aquest procés suposava una elevada càrrega administrativa i, alhora, una alta taxa d'errors durant la generació d'informes de despeses. Per tal de resoldre aquest problema, s'ha dissenyat i desenvolupat una aplicació web \textit{full-stack} que centralitza tot el flux de treball en una única plataforma digital i automatitza l'extracció d'informació dels tiquets i rebuts de despesa.

Aquesta aplicació s'ha desenvolupat utilitzant .NET 8 per a la capa del backend, React amb TypeScript per a la capa del frontend i Microsoft Azure per a la provisió de recursos al núvol. L'arquitectura del sistema segueix els principis de \textit{Clean Architecture} i \textit{SOLID} amb l'objectiu d'obtenir una aplicació mantenible, escalable i fàcilment integrable amb serveis externs. 

A més, l'arquitectura combina processos síncrons i asíncrons. Gràcies a aquesta estratègia, l'aplicació ofereix una experiència d'usuari ràpida i fluida, mentre que les tasques que requereixen un elevat cost computacional, com l'extracció de dades mitjançant intel·ligència artificial, s'executen en segon pla.

Amb la finalitat de garantir l'èxit del projecte, els requisits de l'aplicació s'han prioritzat mitjançant la tècnica MoSCoW, fet que ha permès implementar els requisits de major prioritat dins del temps i del pressupost establerts. Els resultats obtinguts demostren que la solució resol eficaçment el problema identificat, assolint una millora mitjana de l'eficiència del 86,5 \%.
\end{abstract}

\renewcommand{\abstractname}{Abstract}